\section{Discussion and limitations}
The strongest scientific result in this pilot is the measured fidelity
of selected local DFT labels. The cost reduction is useful only because
the retained core forces are compared directly with full-parent
references. The failed two-hop setting illustrates why a pooled average
alone is insufficient. The C32 transfer further shows that a setting
chosen on a smaller molecule can remain useful on a larger parent,
although its near-threshold maximum error calls for caution.

The strongest systems result is that those labels can pass through
the complete dynamics--monitoring--reference--training cycle with
recorded provenance, weighted supervision, promotion, and recovery.
This integration complements prior local active-learning methods;
it does not establish priority for local selection or closed-loop
learning. The platform's physics-aware features are candidate
diagnostics. Their value must be measured against reference-force
errors and simpler acquisition rules, rather than inferred from
their physical motivation.

Several observations limit that interpretation here. The static
comparison contains only two repetitions and a small, repeatedly
examined final-evaluation set. Its PHAL thresholds trigger every
available static frame, and the no-LJ comparison does not reveal
a consistent advantage from the LJ feature. An earlier reference-error
audit labels all 64 evaluated regions positive at each of the
0.1, 0.2, and 0.4~\forceunit\ thresholds. The audit accordingly
does not report average precision as an informative discrimination
metric. With no negative regions, the experiment cannot assess
specificity in distinguishing acceptable from erroneous environments.
Event buffering also has no demonstrated benefit from the static-pool
comparison, which scores unrelated configurations independently.

There is a further distinction between electronic locality and model
locality. A completed diagnostic on one central carbon of one C16
structure finds trained-model force sensitivity extending to roughly
7.6~\AA; truncation to 3~\AA\ can alter that force, whereas
truncation at 8~\AA\ falls below the diagnostic's
$10^{-4}$~\forceunit\ tolerance. The random model has a much smaller
effective range in the same probe. These observations show that a
single-layer cutoff or an untrained-model probe does not determine
the context needed by a trained force prediction. The selected
fragment's DFT-fidelity pass therefore does not eliminate a possible
parent--fragment representation mismatch. This diagnostic is limited
to one model, structure, atom, and perturbation procedure.

The closed-loop topology improvement is encouraging but its initial
model is not a sound baseline for claiming extrapolation repair.
Nor can 2~ps trajectories establish equilibrium conformational
statistics, rare-event kinetics, or long-time stability. An energy
jump of 2.246~eV occurs at a fixed-coordinate model replacement;
this is a change of potential-energy function, not integrator drift.
Energy conservation across model switches is therefore not claimed.
The periodic-box checks support the short trajectories actually
run, not arbitrary later molecular extension.

The evidence supports a platform and an exploratory local-label study.
A stronger repair claim requires an initially qualified force field,
independent confirmation of out-of-domain failures, paired controls
for extra training and acquisition, and a separate final assessment
of both repair and retention of in-domain accuracy. No outcomes
from such an unfinished experiment are included here.
